\documentclass[11pt]{letter}
\usepackage[margin=1in]{geometry}
\usepackage{hyperref}
\signature{Jin Ho Choi}
\address{SmartEAR AI \\ Daegu, South Korea \\ jinhochoi@smartear.co.kr}
\begin{document}
\begin{letter}{Editor-in-Chief \\ IEEE Transactions on Aerospace and Electronic Systems}
\opening{Dear Editor-in-Chief,}

We submit our manuscript, ``A Closed-Loop Subspace--RFS Framework for
Multi-Target DOA Tracking,'' for consideration as a regular paper in the IEEE
Transactions on Aerospace and Electronic Systems.

The paper addresses a core TAES problem---multi-target direction-of-arrival
(DOA) tracking with compact sensor arrays---and its contributions sit squarely
in the Transactions' scope of target tracking, random-finite-set (RFS)
filtering, and radar systems:

\begin{itemize}
\item \textbf{A gated closed-loop architecture} that replaces the standard
  open-loop estimator-to-tracker cascade: the RFS tracker's motion-compensated
  prediction is fed back through a two-stage validation gate to refine the DOA
  estimator's signal subspace. The gate guarantees exact fallback to the
  open-loop estimate, so the deployed system never underperforms the cascade it
  replaces; at low SNR it lowers localization error by up to $56\%$, and under
  target motion it cuts track identity switches by $15\%$.
\item \textbf{Underdetermined tracking capability}: with a higher-order-cumulant
  front-end (COP), a compact $M{=}8$ array tracks $K{=}10>M{-}1$ targets,
  detecting significantly more targets than a MUSIC--PHD cascade. The back-end
  is interchangeable across track-oriented PHD, labeled multi-Bernoulli (LMB),
  and $\delta$-GLMB filters, and a physics-based constant-acceleration motion
  model makes the $\delta$-GLMB back-end best under maneuver.
\item \textbf{Supporting guarantees}: a posterior Cram\'er--Rao bound for the
  complete closed loop (identifiability is restored exactly where the data-only
  bound diverges) and a minimum-variance subspace-fusion proof with a
  gate-enforced benefit condition.
\end{itemize}

The evaluation is Monte-Carlo throughout (up to $500$ trials with $95\%$
confidence intervals), against MUSIC, closed-loop MUSIC, LMB, and $\delta$-GLMB
baselines, and we explicitly delineate the operating regime of each ingredient
rather than claiming uniform superiority.

This manuscript is original, has not been published, and is not under
consideration by any other journal. Thank you for considering our work.

\closing{Sincerely,}
\end{letter}
\end{document}
